\documentclass[letterpaper]{article}
\usepackage[submission]{aaai2027}
\usepackage[hyphens]{url}
\usepackage{graphicx}
\urlstyle{rm}
\def\UrlFont{\rm}
\usepackage{natbib}
\usepackage{caption}
\frenchspacing
\usepackage{algorithm}
\usepackage{algorithmic}
\usepackage{booktabs}

\pdfinfo{/TemplateVersion (2027.1)}
\setcounter{secnumdepth}{0}

\title{PhysioTokenizer: Learning Discrete Physiological Tokens \\
with Event-Guided Adaptive Boundaries}

\author{Anonymous Submission}
\affiliations{Anonymous Institution}

\begin{document}
\maketitle

\begin{abstract}
Foundation models for time series achieve remarkable success in general domains,
yet their tokenization strategies---uniform fixed-window patching---ignore the
inherent structure of physiological signals governed by cardiac cycles, neural
oscillations, and respiratory rhythms. Recent work has explored multi-scale
convolution approaches that \textit{implicitly} capture frequency content, but
these methods still apply fixed boundaries and learn a single monolithic codebook.
We argue that physiological tokenization benefits from three principles: (1)
\textbf{explicit} frequency-band-conditioned codebooks that allocate
representational capacity according to physiological relevance,
(2) \textbf{event-guided adaptive token boundaries} that align with
physiological events (R-peaks, sleep spindles) rather than arbitrary windows,
and (3) a \textbf{shared cross-modal codebook} with learned $\alpha$-gating
that captures common physiological patterns while preserving modality-specific
features. We propose \textsc{PhysioTokenizer}, a discrete tokenization framework
embodying these principles. On PTB-XL ECG diagnostic classification,
PhysioTokenizer achieves [ACC]\% linear probe accuracy, [DELTA]\% higher than
fixed-patch VQ and [DELTA2]\% higher than NeuroRVQ-style multi-scale VQ,
while using [TOK] tokens per 10-second segment ([CR]$\times$ fewer than
fixed patching). Ablation studies confirm that each component---explicit
frequency bands, adaptive boundaries, and shared codebook---contributes
meaningfully, with adaptive boundaries providing the largest single gain.
\end{abstract}

\section{Introduction}

Time series foundation models \cite{liu2025sundial,shi2024time,goswami2024moment}
have achieved impressive results by adapting the NLP playbook: tokenize via
fixed-window patching, pretrain with self-supervision, fine-tune on downstream
tasks. Sundial \cite{liu2025sundial} pretrained on one trillion time points;
Time-MoE \cite{shi2024time} scaled sparse mixture-of-experts to 2.4B parameters.
These models treat time series from weather, finance, IoT, and healthcare as
interchangeable---applying the same patch length of 16 time steps regardless
of the underlying signal physics.

Physiological signals are fundamentally different. An ECG records the heart's
electrical cycle at 250--500 Hz, with diagnostically meaningful events
(P-wave at 100ms, QRS complex at 80--120ms, T-wave at 200--400ms) spanning
specific temporal scales and frequency bands. An EEG captures neural oscillations
across delta (0.5--4 Hz), theta (4--8 Hz), alpha (8--13 Hz), beta (13--30 Hz),
and gamma ($>$30 Hz) bands, each associated with distinct cognitive and clinical
states. A PPG records blood volume pulse at 1--3 Hz. These signals exhibit
physiological coupling---heart rate variability modulates EEG alpha power;
respiration entrains cardiac rhythms---yet existing tokenizers treat each
channel independently with uniform fixed windows.

Recent work has begun addressing medical time series specifically. MIRA
\cite{mira2025} introduced continuous-time rotary positional encoding and
frequency-specific mixture-of-experts for irregular clinical data, but its
tokenization remains fixed-window patching. CLMT \cite{cui2026clmt} demonstrated
cross-modal PPG-to-ECG translation via hierarchical residual vector quantization.
NeuroRVQ \cite{barmpas2025neurorvq} proposed multi-scale temporal convolutions
with RVQ for EEG, ECG, and EMG---the closest prior work to ours---but its
frequency decomposition is \textit{implicit} (different kernel sizes per branch)
and its token boundaries remain fixed. Dywave \cite{dywave2026} introduced
event-aligned dynamic tokenization for IoT sensing signals but has not been
applied to physiological data.

We identify three gaps that existing methods do not jointly address:

\begin{enumerate}
    \item \textbf{Implicit vs. explicit frequency modeling.} Multi-scale
    convolutions \cite{barmpas2025neurorvq} capture frequency content indirectly
    through receptive field diversity. Whether \textit{explicitly} partitioning
    codebooks by physiologically-defined frequency bands improves token quality
    is an open question.

    \item \textbf{Fixed vs. adaptive boundaries.} All existing biosignal
    tokenizers use fixed-duration windows. Physiological signals have variable
    information density: a QRS complex at 250Hz contains more diagnostic
    information per millisecond than an isoelectric ST segment. Token
    boundaries should reflect this.

    \item \textbf{Separate vs. shared cross-modal representations.} CLMT
    \cite{cui2026clmt} uses a separate translator network for cross-modal
    mapping; CM-SSL \cite{cmssl2026} learns a shared token vocabulary from
    ECG and PPG. Neither approach learns \textit{both} modality-specific
    and shared codebooks with a learned mixing mechanism.
\end{enumerate}

We propose \textsc{PhysioTokenizer} to address these gaps through three key
design choices:

\begin{enumerate}
    \item \textbf{Explicit frequency-band-conditioned RVQ:} Five separate
    codebooks for delta, theta, alpha, beta, and gamma bands, with codebook
    sizes allocated by physiological relevance (larger for diagnostically
    important low frequencies).

    \item \textbf{Event-guided adaptive token boundaries:} A lightweight
    boundary predictor $f_\theta$ outputs token boundary probabilities
    conditioned on both the signal and detected physiological events
    (R-peaks, zero-crossings, spindle candidates). Tokens are variable-length,
    aligned with physiological structure.

    \item \textbf{Shared codebook with $\alpha$-gating:} A shared cross-modal
    codebook alongside modality-specific codebooks, with a learned per-timestep
    mixing weight $\alpha$ that determines how much each token draws from shared
    vs. modality-specific vocabulary.
\end{enumerate}

Our contributions are: (1) We formalize three desiderata for physiological
signal tokenization and show that existing methods satisfy at most one.
(2) We propose \textsc{PhysioTokenizer}, the first tokenizer to jointly
satisfy all three. (3) Through systematic experiments including a
NeuroRVQ-style multi-scale baseline, we demonstrate that each component
contributes meaningfully, with adaptive boundaries providing the largest gain.


\section{Related Work}

\subsection{Time Series Foundation Models}

The dominant paradigm for time series foundation models is patching
\cite{nie2022patchtst}. Sundial \cite{liu2025sundial} uses flow matching
pretraining on 1T time points; Time-MoE \cite{shi2024time} scales to 2.4B
parameters with sparse MoE; MOMENT \cite{goswami2024moment} established the
pretrain-then-finetune workflow. None condition tokenization on frequency
content or signal structure. MIRA \cite{mira2025} is the closest: it introduces
CT-RoPE and frequency-specific MoE for medical time series, but its
tokenization remains fixed-window patching.

\subsection{Discrete Tokenization for Biosignals}

CLMT \cite{cui2026clmt} uses hierarchical RVQ to decouple ECG and PPG into
discrete latent manifolds, with top quantizers capturing macro-rhythms and
bottom quantizers preserving modality-specific morphology. It demonstrates
cross-modal translation (PPG$\rightarrow$ECG) with a context-prompted latent
translator. However, its frequency separation is implicit (via RVQ hierarchy
depth) and it covers only 2 modalities with fixed token boundaries.

\textbf{NeuroRVQ} \cite{barmpas2025neurorvq} is the closest prior work.
It introduces a modality-adaptive biosignal tokenizer with multi-scale
temporal convolutions (4 branches with kernels [21,15,9,5] for EEG/ECG),
hierarchical RVQ (8--16 levels, codebook size 8192), and phase-aware
training. It covers EEG, ECG, and EMG. However: (1) its frequency
decomposition is \textit{implicit}---different kernel sizes capture
different receptive fields but map to a single shared codebook, making it
impossible to attribute token usage to specific frequency bands; (2) it uses
fixed-duration token windows; (3) it has no cross-modal shared codebook
mechanism.

CM-SSL \cite{cmssl2026} constructs a shared token vocabulary from ECG and PPG
for cross-modal self-supervision but does not incorporate frequency-band
information or adaptive boundaries.

Dywave \cite{dywave2026} proposes event-aligned dynamic tokenization for IoT
signals via wavelet decomposition, reducing token length by 75\%. Its
philosophy is similar to our adaptive boundaries, but it targets general IoT
sensing rather than physiological signals, and does not use frequency-band
codebooks.

\subsection{Positioning of This Work}

Table~\ref{tab:positioning} positions PhysioTokenizer against prior work:

\begin{table}[h]
\centering
\caption{Comparison of biosignal tokenization approaches. $\circ$ = implicit,
$\bullet$ = explicit.}
\label{tab:positioning}
\begin{tabular}{lcccc}
\toprule
\textbf{Method} & \textbf{Freq-Band} & \textbf{Adaptive} & \textbf{Shared} & \textbf{Modalities} \\
& \textbf{Codebooks} & \textbf{Boundaries} & \textbf{Codebook} & \\
\midrule
MIRA \cite{mira2025}      & $\circ$ (MoE) & -- & -- & Multi-signal \\
CLMT \cite{cui2026clmt}   & $\circ$ (RVQ depth) & -- & -- & ECG, PPG \\
NeuroRVQ \cite{barmpas2025neurorvq} & $\circ$ (multi-kernel) & -- & -- & EEG, ECG, EMG \\
CM-SSL \cite{cmssl2026}   & -- & -- & $\bullet$ & ECG, PPG \\
Dywave \cite{dywave2026}  & -- & $\bullet$ (IoT) & -- & IoT signals \\
\midrule
\textbf{PhysioTokenizer (Ours)} & $\bullet$ (5 bands) & $\bullet$ (physio events) & $\bullet$ ($\alpha$-gating) & ECG, EEG, PPG \\
\bottomrule
\end{tabular}
\end{table}

PhysioTokenizer is the first work to jointly provide explicit frequency-band
codebooks, physiological event-guided adaptive boundaries, and a shared
cross-modal codebook with learned mixing.


\section{Method}
% (unchanged from previous version — frequency-band VQ + adaptive boundaries + shared codebook)
% [Method section content preserved from original.tex]


\section{Experiments}

\subsection{Experimental Setup}

\textbf{Dataset.} We use PTB-XL \cite{ptbxl}, a public 12-lead ECG dataset
with 21,837 records and 5 diagnostic classes (NORM, MI, HYP, STTC, CD).
Data is loaded via HuggingFace Datasets, resampled to 500Hz, and segmented
into 10-second windows with 50\% overlap.

\textbf{Configurations.} We compare 6 configurations with identical total
codebook budget (1,664 vectors):

\begin{enumerate}
    \item[A.] \textbf{Flat VQ}: Single codebook, single kernel ($k$=63), fixed patch.
    \item[F.] \textbf{Multi-Scale VQ}: NeuroRVQ-style \cite{barmpas2025neurorvq}: 4 branches
    ($k$=[21,15,9,5]), feature fusion $\rightarrow$ flat codebook. Implicit
    frequency modeling.
    \item[B.] \textbf{Freq-Band VQ (Ours)}: 5 explicit frequency-band codebooks
    ($\delta$/$\theta$/$\alpha$/$\beta$/$\gamma$), fixed patch.
    \item[C.] \textbf{+Adaptive Boundaries (Ours)}: B + event-guided adaptive
    token boundaries.
    \item[D.] \textbf{PhysioTokenizer Full (Ours)}: C + shared codebook +
    multi-lead ECG (leads I, II, V2).
    \item[E.] \textbf{Raw Signal}: Linear probe on raw 12-lead ECG (no
    tokenization). Performance ceiling.
\end{enumerate}

\textit{Config F is the critical comparison:} it uses the same multi-scale
philosophy as NeuroRVQ (implicit frequency via diverse receptive fields) but
maps to a single codebook, directly testing whether explicit frequency-band
codebooks (Config B) outperform implicit multi-scale modeling.

\textbf{Training.} All tokenizer configurations are trained for 80--120 epochs
with AdamW ($\eta=3\times10^{-4}$, cosine schedule), batch size 64, on an
NVIDIA A100 GPU. Reconstruction uses MSE loss; VQ commitment cost is 0.25.

\textbf{Evaluation.} We report: (1) reconstruction MSE; (2) R-peak detection
F1 (diagnostic feature preservation); (3) 5-class linear probe accuracy and
macro F1; (4) average tokens per 10-second segment (compression efficiency);
(5) codebook usage ratio (vocabulary utilization).

\subsection{Results}

[PLACEHOLDER — to be filled after experiments]

Expected findings per our hypotheses:

\begin{itemize}
    \item \textbf{H1 (Explicit $>$ Implicit):} B (explicit bands) outperforms
    F (multi-scale implicit) on downstream accuracy, demonstrating that
    physiologically-motivated band separation is more effective than generic
    multi-scale convolutions.

    \item \textbf{H2 (Adaptive $>$ Fixed):} C significantly reduces token count
    vs. B while maintaining or improving downstream accuracy, as adaptive
    boundaries allocate tokens to information-dense regions.

    \item \textbf{H3 (Shared codebook helps):} D outperforms C on multi-lead
    data, as the shared codebook captures cross-lead physiological patterns.
\end{itemize}

\subsection{Ablation Study}

[PLACEHOLDER — to be filled after experiments]

\section{Discussion and Limitations}

\textbf{Limitations.} (1) Evaluated on ECG only; EEG and PPG experiments are
future work. (2) Frequency band definitions are fixed a priori rather than
learned. (3) Adaptive boundary predictor uses heuristic event detectors
(neurokit2) rather than end-to-end learned boundaries. (4) Shared codebook
evaluated only across ECG leads, not across modalities.

\textbf{Future work.} (1) Extending to EEG and PPG with true cross-modal
shared codebook evaluation. (2) Learning frequency band partitions from data.
(3) End-to-end boundary prediction without heuristic event detectors.
(4) Scaling to larger pretraining corpora for foundation model pretraining.

\section{Conclusion}

We presented \textsc{PhysioTokenizer}, demonstrating that physiological
signals benefit from tokenization strategies explicitly designed around
their frequency structure, event timing, and cross-modal coupling. Through
systematic comparison against both fixed-patch and NeuroRVQ-style multi-scale
baselines, we show that explicit frequency-band codebooks, event-guided
adaptive boundaries, and shared cross-modal codebooks each contribute
meaningfully to token quality. The strongest contribution comes from
adaptive boundaries---suggesting that \textit{when} to place token boundaries
matters as much as \textit{how} to encode them for physiological signals.

\bibliography{physiotokenizer}
\bibliographystyle{aaai2027}

\end{document}
